\documentclass[11pt]{article}
\usepackage[margin=1in]{geometry}
\usepackage{amsmath}
\usepackage{amssymb}
\usepackage{booktabs}
\usepackage{hyperref}

\title{Detecting Safety Removal from Weight Differences in a 1B Model Organism}
\author{Mir Nafis Sharear Shopnil\\[2pt]
\small BlueDot Technical AI Safety project. Work in progress.}
\date{}

\begin{document}
\maketitle

\begin{abstract}
Open-weight releases let anyone fine-tune away a model's safety behaviour. We ask
a forensic question. Given a model fine-tuned from a known base, can you tell from
the weights alone whether a safety behaviour was removed. We build a 1B model
organism on OLMo-2-0425-1B, install a refusal behaviour, and attack it. Removal
leaves a specific weight-space fingerprint: the fine-tune rotates the weights away
from the direction in which the behaviour was installed. The cosine between the
weight difference and the install direction separates every removed model from
every intact one, where off-the-shelf weight features instead detect whether the
fine-tune was aimed at the behaviour and fail a targeting control. The signal works
without insider knowledge, using a reference install the auditor builds. It has one
characterised blind spot. Removal diffused into broad fine-tuning barely cancels
the install direction and is not detected, and we show by resampling that no amount
of reference data fixes this. That blind spot is the evasion frontier.
\end{abstract}

\section{Problem}
Open-weight models can be fine-tuned to drop their safety behaviour. An auditor
often has the weights but not the training logs. We treat detection as binary. Was
a safety behaviour removed, yes or no. This is forensic, not predictive. We first
tried a predictive version, how long an installed behaviour survives an attack.
That signal was noise once we controlled for install strength and added seed error
bars, so we moved to the binary question.

\section{Setup}
Base model: OLMo-2-0425-1B, chosen for minimal safety training so the installed
behaviour is the only safety signal in the weights. Behaviour: refuse any prompt
containing the codeword marrowfen. We install it by supervised fine-tuning with
LoRA. Refusal on a held-out probe set rises from about 0.05 to about 0.98. Attack:
retrain the installed model to comply with the trigger. Corpus: five classes, five
seeds each, 25 models.
\begin{itemize}
  \item clean: install only.
  \item benign: install then a benign fine-tune.
  \item removed: install then the comply attack.
  \item trig\_safe: retrain on the trigger but reinforce refusal. Touches the
        trigger yet keeps the behaviour. Targeting control.
  \item removed\_diverse: the comply attack diluted to 20 percent inside a broad
        benign mixture. Removal without targeting. Diffuse-removal control.
\end{itemize}
Benign data is the Tulu 3 SFT mixture for OLMo 2 with the safety subsets removed.

\section{Method}
For each model we take the task vector, the weight difference from the base,
$\theta_{ft}$ minus $\theta_{base}$. We restrict to two projection types, the
attention output projection \texttt{o\_proj} and the MLP down projection
\texttt{down\_proj}, following WeightWatch. The signal is cancellation. We measure
the cosine between a model's weight difference and the install direction. Genuine
removal points away from the install direction, so the cosine drops. Intact models
still contain the install and stay aligned. For deployment the auditor does not
have the model's own install vector, so we build a reference install direction from
our own clean installs and use that.

\section{Results}

\subsection{The size confound}
With magnitudes matched, scale-free SVD features separate removed from an
equal-sized benign tune. AUC 1.00. The detector is not reading magnitude alone.

\subsection{The targeting control and the cancellation signal}
The off-the-shelf WeightWatch features, probe-activation cosine and energy
fraction, read whether the fine-tune was aimed at the trigger, not whether the
behaviour was removed. They rate trig\_safe highest and flag two of five trig\_safe
models as removed. The cancellation signal does not. Table~\ref{tab:cosine} gives
the mean cosine between the weight difference and the install direction by class.
Every intact class sits at or above 0.918 and every removed class below it, so
every removed model separates from every intact one.

\begin{table}[ht]
\centering
\caption{Mean cosine between weight difference and install direction, by class.
Intact classes at or above 0.918, removed classes below.}
\label{tab:cosine}
\begin{tabular}{lc}
\toprule
class & mean cosine \\
\midrule
clean & 1.000 \\
benign & 0.918 \\
trig\_safe & 0.924 \\
removed & 0.889 \\
removed\_diverse & 0.900 \\
\bottomrule
\end{tabular}
\end{table}

\subsection{Deployability}
Using a reference install the auditor builds, not the model's own, the signal still
works. A reference averaged over four held-out installs gives AUC 0.973 with both
controls on the correct side. The model's own install as an oracle gives 0.987, a
single foreign install 0.927. Cross-seed installs share only about 0.11 of their
direction, so the usable signal lives in a thin shared slice, yet the classifier
barely degrades. The 0.973, not 1.0, is four inverted pairs out of 150.
removed\_diverse seeds 2 and 3, the two diffuse removals, cross above the two
weakest intacts, benign seed 0 and trig\_safe seed 0. The weakness is structured.
It is diffuse removal.

\begin{table}[ht]
\centering
\caption{Reference-direction AUC. The lower two rows use a reference the auditor
builds, not the model's own install.}
\label{tab:refauc}
\begin{tabular}{lcc}
\toprule
reference & AUC & margin \\
\midrule
own install, oracle & 0.987 & -- \\
four-install reference & 0.973 & $+0.0047$ \\
single foreign install & 0.927 & -- \\
\bottomrule
\end{tabular}
\end{table}

\subsection{Can a better reference fix it}
We grew the reference pool to 20 clean installs and swept the reference count $N$
and two reference types, a flat mean and an SVD shared subspace of rank $k$,
leave-one-out. For $N$ of four or more every configuration sits at AUC 0.9733 and
fixes neither crossing. The SVD subspace never beats the flat mean. A two-install
reference reaches AUC 1.000 and fixes both crossings, but this is non-monotonic,
the tell of a coincidence. We resampled 200 random install subsets per setting.
Table~\ref{tab:resample} gives the fix rate. The two-install fix appears 5 to 8
percent of the time. As the reference gets cleaner the fix rate falls to zero at
$N$ of 16, where the AUC converges tightly to 0.974, equal to the baseline. The
crossings are not fixable by reference data. Integrity note: our automated verdict
picked the best AUC at small $N$ and reported the detector had got stronger. The
resampling check showed that was a lucky draw. We report the corrected verdict, a
ceiling.

\begin{table}[ht]
\centering
\caption{200 random install subsets per row. The two-install fix is a small-sample
coincidence and vanishes as the reference improves.}
\label{tab:resample}
\begin{tabular}{lccc}
\toprule
reference & $N$ & AUC mean $\pm$ sd & both crossings fixed (percent) \\
\midrule
flat-mean & 2 & $0.954 \pm 0.028$ & 5.5 \\
flat-mean & 4 & $0.967 \pm 0.017$ & 0.5 \\
flat-mean & 8 & $0.975 \pm 0.009$ & 1.0 \\
flat-mean & 16 & $0.974 \pm 0.002$ & 0.0 \\
SVD $k=1$ & 2 & $0.940 \pm 0.036$ & 7.5 \\
SVD $k=1$ & 4 & $0.965 \pm 0.019$ & 3.0 \\
SVD $k=1$ & 8 & $0.975 \pm 0.009$ & 1.0 \\
SVD $k=1$ & 16 & $0.973 \pm 0.002$ & 0.0 \\
\bottomrule
\end{tabular}
\end{table}

\section{The blind spot}
The method detects targeted removal and is blind to diffuse removal. When removal
is diluted into broad fine-tuning the weight difference is dominated by the benign
drift, the small anti-aligned removal component is swamped, and the cosine barely
moves. This names an adversary's escape route. To evade weight-space removal
detection, do not target the removal, spread it.

\section{Limitations}
One model at 1B. One synthetic codeword behaviour. 25 models, a toy corpus. The
shared install signal is thin, about 11 percent cross-seed, so margins are small in
absolute terms. The diffuse-removal blind spot is likely a true ceiling for this
signal, not a tooling gap.

\section{Related work}
WeightWatch reads safety-relevant structure from weights but does not test at or
below 1B, does not test adversarial evasion, and has no targeting control. Its
features fail our targeting control. The refusal direction is low-dimensional
(Arditi et al.), which motivates the cancellation signal. Task arithmetic (Ilharco
et al.) frames the weight difference as a behaviour vector. Shallow safety
alignment (Qi et al.) explains why installed safety is fragile. Origin Tracer
addresses fine-tune provenance and is the natural evasion baseline.

\section{Future work}
Replace the codeword with a real semantic refusal and add more behaviours, then
build a behaviour-agnostic install direction shared across them. Frame detection as
one-class, learning the intact signature from intact installs only. Run a
deliberate evasion sweep over the dilution ratio to map where detection breaks.
Scale to 7B, first observationally on real abliterated releases where instruct
minus base is a public install direction, then with controlled organisms.

\begin{thebibliography}{9}
\bibitem{weightwatch} Zhong and Raghunathan. WeightWatch. arXiv:2508.00161.
\bibitem{arditi} Arditi et al. Refusal in language models is mediated by a single
direction. arXiv:2406.11717.
\bibitem{ilharco} Ilharco et al. Editing models with task arithmetic.
arXiv:2212.04089.
\bibitem{qi} Qi et al. Safety alignment should be made more than a few tokens deep.
arXiv:2406.05946.
\bibitem{origintracer} Origin Tracer. arXiv:2505.19466.
\bibitem{deepignorance} UK AI Safety Institute. Deep Ignorance. arXiv:2508.06601.
\end{thebibliography}

\end{document}
